\section{Weighted minimization}

In previously presented methods, it is assumed that all measurements are equally important. However, this may not be the case. As an example, consider that some of the measurements were taken when the body of water was being disturbed by heavy water currents, affecting the position of the sensors and the perceived sound speed in water. These measurements are useful, but their data is knowingly less accurate and reliable than that of other measurements, and therefore a weighting step for these measurements can be taken.

Mathematically, this is modeled by adapting \cref{subeq:sec2:iterative_minimization:eq1}. This equation can be expanded, and written as
\begin{equation}
	\bvz{i+1} = \arg \min_{\bvz} \sum_{i=1}^{M} \pts{\bvj[T]{m;i} \bvz - t_{m;\obs}}^2
\end{equation}
in which $\bvj{m;i}$ is the distance traversed in each $i$-th cell, by the $m$-th ray, and $t_{m;\obs}$ is the observed total travel-time for the $m$-th ray. With this formulation, a weighting step can be easily added, resulting in
\begin{equation}
	\bvz{i+1} = \arg \min_{\bvz} \sum_{i=1}^{M} w_i \pts{\bvj[T]{m;i} \bvz - t_{m;\obs}}^2
\end{equation}
where $w_i > 0$ is the weight for the $i$-th measurement. Returning to the matricial formulation, we get
\begin{equation}
	\bvz{i+1} = \arg \min_{\bvz} \norm{\bvW \pts{\bvJ{i} \bvz - \bvt{\obs}}}^2
\end{equation}
where $\bvW$ is a diagonal matrix, such that the $\el{\bvW}[i,i] = \Sqrt{w_i}$. Using the SVD and pseudo-inverse approach as established, the solution to this minimization is
\begin{equation}
	\bvG{i} = \pinv*{\bvW \bvJ{i}} \bvW
\end{equation}

